% =====================================================================
%  TITRE I : DE L'ORGANISATION DU SERVICE
% =====================================================================
\titrepartie{Titre I}{De l'organisation du service}

% ---------------------------------------------------------------------
\article{Principes fondateurs}

\chapeau{Le Los Santos Police Department est un service public. Le service rendu aux citoyens doit être irréprochable. Le présent article pose les principes qui encadrent les relations entre les personnels, avec les autres services publics et avec la population. Un manquement à cet article peut, selon sa gravité, conduire jusqu'à la révocation.}

\cln[respect]{Respect}{Chaque agent est tenu au respect envers ses collègues, quel que soit leur grade, envers les personnels des autres services publics et envers les citoyens. Cette obligation est inconditionnelle : elle ne dépend ni du comportement de l'interlocuteur, ni du contexte de l'échange. L'agent représente en toutes circonstances l'image du département.}

\cln[camaraderie]{Camaraderie}{Les agents se doivent assistance mutuelle. Laisser sciemment un collègue seul dans une situation dangereuse ou dégradée constitue une faute grave, passible d'une mise à pied pouvant aller jusqu'à la révocation sans préavis.}

\cln{Coopération}{Chaque agent coopère pleinement avec les autres services publics dès lors que la situation l'exige. Cette coopération s'exerce sans réserve, sans rétention d'information et sans considération de rivalité entre services.}

\cln[professionnalisme]{Professionnalisme}{L'agent demeure professionnel en toutes circonstances. Les débordements comportementaux, les propos déplacés, la familiarité excessive avec le public et toute attitude de nature à discréditer le département sont prohibés.}

\cl{Entre eux, les agents s'appellent par leur nom de famille ou par leur grade. L'usage du prénom est réservé aux échanges hors service ou hors présence du public.}

\cl[badge]{Tout agent est tenu de communiquer son numéro de badge lorsqu'il lui est demandé, que la demande émane d'un citoyen, d'un avocat ou d'un personnel d'un autre service. Le refus de communiquer son numéro de badge constitue une faute disciplinaire caractérisée.}

\cln[harcelement]{Harcèlement et discrimination}{Tout comportement de harcèlement, toute conduite déplacée et tout propos discriminatoire fondés notamment sur le sexe, l'origine, l'orientation, l'apparence ou les croyances exposent leur auteur à une révocation sans préavis.}

\cln[insubordination]{Insubordination}{L'insubordination est strictement interdite. Chaque agent respecte l'autorité de ses supérieurs hiérarchiques et exécute les ordres reçus dans le cadre du service. Refuser d'exécuter un ordre légitime, contester publiquement l'autorité d'un supérieur ou adopter une attitude de défiance envers la hiérarchie constitue une violation grave de la discipline.}

\clsuite{Un agent peut et doit refuser d'exécuter un ordre \emph{manifestement illégal} ou dont l'exécution mettrait gravement en danger sa vie, celle d'un collègue ou celle d'un tiers. Ce refus doit être exprimé immédiatement, motivé, et porté sans délai à la connaissance d'un superviseur d'un échelon supérieur ainsi que, le cas échéant, de la Division des Affaires Internes. Le refus légitime d'un ordre illégal ne peut donner lieu à aucune sanction ni à aucune mesure de rétorsion.}

\cln[confidentialite]{Confidentialité}{Les informations obtenues dans l'exercice des fonctions sont confidentielles et ne peuvent être divulguées à des personnes non habilitées, sauf nécessité légale. Cette obligation vise en particulier les enquêtes en cours, l'identité des témoins et des suspects, les antécédents, les adresses et les stratégies opérationnelles. Sa violation expose à des sanctions sévères pouvant aller jusqu'à la révocation et à la saisine de l'autorité judiciaire.}

\cln[conflits]{Règlement des différends}{Tout désaccord entre agents se règle par la voie hiérarchique, dans l'ordre suivant : Senior Lead Officer, puis Sergeant, puis Command Staff et, en dernier recours, État-Major. Un différend ne doit en aucun cas être exposé sur les ondes, devant le public ou devant les personnels d'un autre service, ni perturber l'efficacité du service.}

\cl{Tout agent responsable d'un agent en période probatoire supervise l'intégralité de la procédure engagée par ce dernier et en répond devant la hiérarchie.}

\cl[procedure-terminee]{Tout agent engageant une procédure, seul ou avec un partenaire, est tenu de la mener à son terme. En cas d'impossibilité, il assure la transmission complète du dossier à un collègue et en informe le Watch Commander. L'abandon d'une procédure sans motif légitime ni passation expose à des sanctions disciplinaires.}

% ---------------------------------------------------------------------
\article{Obligations de service}

\cl[quota]{Chaque agent accomplit un minimum de \textbf{sept heures de service par semaine}. L'agent dans l'impossibilité de respecter cette obligation en informe sa hiérarchie et dépose une demande d'absence dans les conditions prévues à l'article relatif à la disponibilité et aux absences.}

\cl{Aucun agent ne peut demeurer plus de \textbf{soixante-douze heures} sans se présenter au service sans avoir préalablement déclaré une absence. Le manquement à cette obligation expose à une procédure de radiation.}

\cl{La prise et la fin de service font l'objet d'une annonce radio et d'une inscription au registre du poste. Aucun agent ne peut exercer les prérogatives de police sans avoir régulièrement pris son service.}

\cl{L'agent se présente à son service en tenue conforme, avec un équipement complet et opérationnel, et dans un état physique et psychologique compatible avec l'exercice de ses fonctions. Il est interdit de prendre son service sous l'emprise de l'alcool ou de toute substance altérant le discernement.}

\cl{Le service prime sur les activités de division. Un agent membre d'une division peut être rappelé en patrouille par le Watch Commander lorsque l'effectif l'exige.}

% ---------------------------------------------------------------------
\article{Protection, bien-être et aptitude des agents}

\chapeau{Le département attache une importance déterminante à la sécurité et au bien-être de ses agents. Un service efficace suppose des agents protégés, soutenus et aptes.}

\cln[protection]{Protection contre le harcèlement et les agressions}{Tout agent confronté à des faits de harcèlement, de menaces ou d'agression, qu'ils émanent d'un supérieur, d'un collègue ou d'un civil, est encouragé à les signaler sans crainte de représailles. Une enquête confidentielle est ouverte dans les meilleurs délais afin de garantir la sécurité et la dignité de l'agent concerné.}

\cl{Toute mesure de rétorsion, directe ou indirecte, prise à l'encontre d'un agent en raison d'un signalement effectué de bonne foi constitue une faute lourde, sanctionnée indépendamment du sort réservé au signalement lui-même.}

\cln[soutien]{Soutien psychologique}{Le département reconnaît les effets possibles du stress, de l'épuisement professionnel et des difficultés psychologiques liées à l'exercice de la fonction. Tout agent peut solliciter un entretien confidentiel auprès d'un membre du Command Staff ou d'un référent désigné de la Management Relations Division. Cette démarche n'entraîne par elle-même aucune sanction ni mesure disciplinaire.}

\cl{À la suite d'un événement grave (usage de la force létale, décès d'un collègue, prise d'otages, blessure en service), l'agent concerné bénéficie d'un temps de retrait et d'un entretien de suivi avant tout retour sur le terrain. Ce retrait est une mesure de protection et ne préjuge d'aucune faute.}

\cln[aptitude]{Aptitude au service}{Tout agent doit se trouver, à sa prise de service, dans un état physique et psychologique compatible avec l'exercice de ses fonctions et le port de l'arme de service. Il appartient à chacun d'y veiller et, en cas de doute, de se signaler à sa hiérarchie plutôt que de prendre son service.}

\cl{Lorsque des éléments objectifs le justifient, le Command Staff peut retirer temporairement à un agent l'autorisation de porter l'arme de service et l'affecter à des fonctions sans port d'arme. Cette mesure est motivée, notifiée par écrit et réexaminée régulièrement. Elle constitue une mesure de protection et non une sanction.}

% ---------------------------------------------------------------------
\article{Chaîne hiérarchique et grades}

\chapeau{Le Los Santos Police Department s'organise en cinq corps. Le respect de la chaîne hiérarchique est impératif et chacun de ses échelons doit être observé.}

\cl[chaine]{La chaîne hiérarchique du département comprend, par ordre décroissant d'autorité : l'État-Major, le Command Staff, la Supervision, les officiers de police et les personnels administratifs. Les grades de Detective constituent une filière d'investigation dont les échelons s'alignent sur ceux de la filière de terrain.}

\sstitre{État-Major}

\cln{Chief}{Il dirige l'ensemble du département et représente l'autorité la plus élevée du service. Il définit les orientations générales de l'organisation, supervise l'intégralité des divisions et veille au maintien de l'ordre, de la discipline et de l'efficacité globale du service. Il assure les relations avec les autres services publics et représente officiellement le département lors des événements majeurs, des réunions institutionnelles et des situations de crise.}

\cln{Commander}{Il assure la coordination stratégique et administrative de son secteur, veille au fonctionnement des unités placées sous son autorité et s'assure de l'application des directives de l'État-Major. Il supervise les opérations d'ampleur, encadre la supervision et participe au développement du service. Membre de l'État-Major, il contribue à la gestion du département et aux décisions stratégiques.}

\sstitre{Command Staff}

\cln{Captain}{Il est responsable d'un poste de police et de l'effectif qui y est affecté. Il répond de la discipline, de l'organisation du service et de la qualité des procédures produites par son poste.}

\cln{Lieutenant}{Il assiste le Captain dans ses missions et constitue l'interlocuteur privilégié de la supervision. Il conduit les revues de procédure et instruit les demandes de sanction transmises par les superviseurs.}

\cl{L'ensemble des membres du Command Staff et de l'État-Major supervise le service dans sa globalité et s'engage à prendre en permanence les mesures propres à l'améliorer.}

\sstitre{Supervision}

\cln{Sergeant II}{Il est à la tête de la supervision. Il organise les patrouilles et en assure le bon déroulement, et constitue le lien direct entre les unités de terrain, le Command Staff et l'État-Major. Il exerce en principe la fonction de Watch Commander.}

\cln{Sergeant I}{Il assure la supervision des unités de patrouille et veille au respect du présent règlement sur le terrain.}

\cl{Les superviseurs contrôlent le travail des unités de patrouille, les organisent sur le terrain et prennent le commandement lors des événements majeurs. Ils peuvent proposer un agent à la promotion et engager les sanctions relevant de leur échelon.}

\sstitre{Officiers de police}

\cln{Police Senior Lead Officer}{Officiers expérimentés s'étant distingués au cours de leur carrière, ou ayant fait le choix de ne pas rejoindre la supervision. Ils assurent l'encadrement des agents probatoires et conduisent les évaluations de First Lincoln.}

\cln{Police Officer III}{Officiers confirmés, présents au sein du service depuis une durée significative. Avec les Senior Lead Officers, ils ont la charge de la formation des nouveaux agents.}

\cln{Police Officer II}{Agents ayant acquis des bases solides et appelés à accompagner les agents du grade inférieur dans l'acquisition de nouvelles compétences.}

\cln{Police Officer I}{Grade d'entrée dans le service. Il comporte deux états successifs, qui ne constituent pas deux grades distincts mais deux étapes d'une même progression.}

\clsuite{\stress{Police Officer I probatoire.} État d'affectation de tout nouvel agent intégrant le service. L'agent est en période d'essai : il exerce sous encadrement permanent, ne patrouille jamais seul et peut être radié à tout moment par le Command Staff ou l'État-Major, sans que cette décision constitue une sanction disciplinaire.}

\clsuite{\stress{Police Officer I confirmé.} État acquis après réussite du First Lincoln. L'agent est apte à patrouiller sans encadrement. En l'absence de supérieur disponible, il peut prendre en patrouille un agent probatoire.}

\sstitre{Personnels administratifs}

\cln{Secrétaire}{Assiste le département dans ses tâches administratives internes. Le secrétariat est responsable du traitement des plaintes, du recensement des civils et de la gestion documentaire du département. Il peut être rattaché à la Management Relations Division pour la logistique : gestion des stocks, commandes internes, prise de rendez-vous et entretiens de suivi des agents probatoires. Il ne participe pas aux interventions mais demeure essentiel au fonctionnement du service.}

\sstitre{Filière d'investigation}

\cl{Le département reconnaît trois échelons de Detective, correspondant à des fonctions d'investigation spécialisées et s'alignant sur la hiérarchie des grades de terrain.}

\vspace{4pt}
\noindent
\begin{tabular}{P{30mm} P{28mm} P{86mm}}
\headrow
\thead{Grade} & \thead{Équivalence} & \thead{Attributions} \\
\tkey{Detective I} & \tcell{Senior Lead Officer} & \tcell{Conduit des investigations approfondies sans assumer de responsabilité de supervision directe.} \\
\altrow
\tkey{Detective II} & \tcell{Sergeant I} & \tcell{Prend en charge les enquêtes complexes, coordonne de petites équipes d'investigation et appuie les superviseurs sur le terrain.} \\
\tkey{Detective III} & \tcell{Sergeant II} & \tcell{Enquêteur senior. Supervise les enquêtes majeures, gère les ressources d'investigation et assure la liaison entre les équipes de terrain et le Command Staff.} \\
\bottomrule
\end{tabular}

\begin{note}[Autorité fonctionnelle]
L'équivalence de grade détermine l'autorité hiérarchique générale. Sur une scène en cours, c'est toutefois le Operation Lead qui commande, quelle que soit la filière d'appartenance des agents présents. Un Detective ne prend pas le commandement d'une intervention de voie publique du seul fait de son grade.
\end{note}

% ---------------------------------------------------------------------
\article{Recrutement, académie et formation terrain}

\chapeau{L'intégration au département suit un parcours unique et non négociable. Aucune étape ne peut être supprimée, quelle que soit l'expérience antérieure du candidat.}

\sstitre{Conditions de recrutement}

\cl[conditions]{Le candidat doit remplir cumulativement les conditions suivantes : casier judiciaire vierge, permis de conduire valide dans l'État de San Andreas, et absence de lien avéré avec une organisation criminelle.}

\cl{Ces conditions sont des conditions d'emploi permanentes : leur perte en cours de carrière entraîne la suspension immédiate des prérogatives de police et l'ouverture d'une procédure.}

\sstitre{Parcours d'intégration}

\begin{enumerate}
  \item \stress{Dossier de candidature} : identité, antécédents, motivation et disponibilité.
  \item \stress{Enquête de moralité} : vérification du casier, des antécédents et des affiliations.
  \item \stress{Entretien de recrutement} : conduit devant un jury composé d'au moins un membre du Command Staff ou de la supervision.
  \item \stress{Académie} : formation initiale dispensée par la Los Santos Police Academy.
  \item \stress{Assermentation} : prestation du serment, remise de la plaque et de l'arme de service.
  \item \stress{Période probatoire} : affectation en Police Officer I probatoire, service encadré, sanctionné par le First Lincoln.
\end{enumerate}

\sstitre{Période probatoire et formation terrain}

\cl{La formation initiale dispensée par l'académie porte sur les fondamentaux du service et sur le présent règlement. L'ensemble des procédures opérationnelles est acquis sur le terrain, pendant la période probatoire, sous la responsabilité du formateur désigné.}

\cl[fto]{Tout agent nouvellement assermenté est affecté en Police Officer I probatoire et placé sous la responsabilité d'un formateur de terrain, Police Officer III ou Senior Lead Officer désigné. Le formateur répond des erreurs de son agent probatoire pendant les deux premières phases.}

\vspace{4pt}
\noindent
\begin{tabular}{P{26mm} P{34mm} P{84mm}}
\headrow
\thead{Phase} & \thead{Intitulé} & \thead{Rôle de l'agent probatoire} \\
\tkey{Phase 1} & \tcell{Observation} & \tcell{L'agent probatoire observe. Il ne conduit pas et ne tient pas la radio.} \\
\altrow
\tkey{Phase 2} & \tcell{Assistance} & \tcell{L'agent probatoire tient la radio et le terminal embarqué. Le formateur conduit et mène les contacts.} \\
\tkey{Phase 3} & \tcell{Conduite} & \tcell{L'agent probatoire mène les interventions et rédige les procédures. Le formateur corrige en fin d'intervention.} \\
\altrow
\tkey{Phase 4} & \tcell{First Lincoln} & \tcell{Évaluation finale. Le formateur n'intervient plus, sauf danger. L'agent probatoire gère seul l'intégralité de la procédure.} \\
\bottomrule
\end{tabular}

\cl[first-lincoln]{Le First Lincoln est conduit par un Senior Lead Officer, un Sergeant ou un membre de la Los Santos Police Academy, distinct autant que possible du formateur habituel. Il donne lieu à une grille d'évaluation écrite et à un avis motivé. Sa réussite fait passer l'agent de Police Officer I probatoire à Police Officer I confirmé.}

\cl{En cas d'échec, la période probatoire est prolongée et un nouveau passage est organisé. Le Command Staff et l'État-Major peuvent mettre fin à la période d'essai à tout moment, notamment après plusieurs échecs ou en cas d'inaptitude manifeste.}

\begin{interdit}[Restrictions applicables au Police Officer I probatoire]
L'agent probatoire ne peut en aucun cas : patrouiller seul ou en unité Lincoln, intégrer une division, se voir délivrer une arme lourde, exercer le commandement d'une intervention, ni encadrer un autre agent. Étant en période d'essai, il peut être radié à tout moment par le Command Staff ou l'État-Major.
\end{interdit}

% ---------------------------------------------------------------------
\article{Avancement, promotions et rétrogradations}

\cl{Les perspectives d'évolution sont ouvertes à tout agent démontrant un investissement sérieux et constant dans ses fonctions. L'ancienneté seule ne fonde aucun droit à l'avancement.}

\cl[promotions]{L'État-Major et le Command Staff sont seuls habilités à prononcer les promotions et les rétrogradations. Ils apprécient les profils au regard des critères définis pour chaque grade et des besoins du service. Un superviseur peut proposer un agent à la promotion ; il ne peut la prononcer.}

\cl{Les critères d'appréciation comprennent notamment : le volume et la régularité du service accompli, la qualité des procédures et des rapports produits, la maîtrise des procédures opérationnelles, le comportement envers le public et les collègues, l'état du casier disciplinaire et la capacité à encadrer.}

\cl{Les cérémonies de promotion se tiennent toutes les deux à quatre semaines selon l'activité du service et les besoins en encadrement.}

\cl{Une rétrogradation est une sanction disciplinaire et obéit aux règles du titre relatif à la discipline. Elle est motivée par écrit et notifiée à l'agent, qui dispose d'une voie de recours.}

% ---------------------------------------------------------------------
\article{Disponibilité, congés et absences}

\cl[quota-rappel]{Tout agent est tenu d'accomplir sept heures de service par semaine et ne peut demeurer plus de soixante-douze heures consécutives sans se présenter au service sans avoir déclaré une absence.}

\cl[loa]{La demande d'absence indique la date de départ et la date de retour envisagée. L'agent peut y joindre un motif ; il est libre de ne pas le faire. Lorsqu'un motif est communiqué, l'agent choisit s'il doit demeurer confidentiel, auquel cas il n'est accessible qu'au Command Staff et à l'État-Major.}

\cl[loa-nombre]{\textbf{Le nombre et la durée des absences ne sont pas limités.} Une absence régulièrement déclarée ne peut en elle-même fonder aucune sanction.}

\cl[loa-frequence]{Le Command Staff et l'État-Major se réservent toutefois le droit de mettre fin aux fonctions d'un agent dont les absences, par leur fréquence ou leur durée cumulée, deviennent incompatibles avec l'exercice de sa fonction, avec son grade ou avec ses responsabilités de division. Cette décision est motivée, notifiée par écrit et ouvre droit à recours.}

\cl{Pendant une absence déclarée, l'agent est dispensé du quota de service mais demeure soumis au présent règlement. Il n'exerce aucune prérogative de police et ne peut se présenter comme agent en fonction.}

\cl{Le retour s'effectue à la date annoncée, sans formalité, avec rétablissement immédiat des obligations de service. Un retour anticipé est possible sur simple information de la hiérarchie.}

\sstitre{Absence non déclarée}

\cl[absence-non-declaree]{L'agent absent sans déclaration préalable fait l'objet d'une relance de sa hiérarchie au bout de \textbf{sept jours}. À défaut de retour ou de réponse dans les \textbf{quatorze jours} suivant le début de l'absence, la radiation est prononcée.}

\cl{Un agent radié pour inactivité peut solliciter sa réintégration dans un délai de trois mois. Elle peut alors être prononcée sans nouveau passage par l'académie, après un entretien et une remise à niveau.}

\cl[demission]{La démission est présentée par écrit à l'État-Major. Elle donne lieu à la restitution intégrale de la plaque, de l'arme de service et de l'équipement, ainsi qu'à un entretien de sortie. La restitution conditionne la délivrance du certificat de fin de service.}

% ---------------------------------------------------------------------
\article{Organisation des patrouilles}

\cl[watchcommander]{Afin d'assurer la couverture du territoire, des agents sont désignés à des postes d'organisation des unités. Le \stress{Watch Commander}, en principe Sergeant II, est le superviseur de référence des unités en patrouille. Il organise les unités et leur supervision pendant la garde, prend le commandement lorsqu'un événement majeur survient et répartit les superviseurs là où ils sont nécessaires.}

\clsuite{Tout personnel patrouillant dans la zone placée sous sa supervision se conforme à ses instructions, indépendamment de son grade ou de son ancienneté.}

\cl{Les patrouilles en binôme sont la règle. Lorsque moins de quatre agents sont en service, les officiers, à l'exception des Police Officer I probatoires, peuvent patrouiller seuls.}

\cl{Aucune division ne peut être déployée si moins de deux patrouilles Adam sont en service. Cette règle ne s'applique pas aux agents du Detective Bureau.}

\cl[communication]{Tout agent utilise, dans ses communications radio, les codes, l'alphabet et les procédures de communication en vigueur. Le non-respect de cette obligation expose à une sanction.}

\sstitre{Types d'unité}

\noindent
\begin{tabular}{P{26mm} P{122mm}}
\headrow
\thead{Indicatif} & \thead{Composition et emploi} \\
\tkey{Lincoln} & \tcell{Unité classique composée d'un agent. Présence sur tous les appels.} \\
\altrow
\tkey{Adam} & \tcell{Unité classique composée de deux agents. Présence sur tous les appels. Format de référence.} \\
\tkey{Tango} & \tcell{Unité classique composée de trois agents. Présence sur tous les appels.} \\
\altrow
\tkey{Queen} & \tcell{Unité classique composée de quatre agents. Présence sur tous les appels.} \\
\tkey{Romeo} & \tcell{Unité d'intervention rapide de la Traffic Division. Présence sur tous les appels.} \\
\altrow
\tkey{Mary} & \tcell{Unité motorisée de la Traffic Division, adaptée aux terrains montagneux.} \\
\tkey{Air} & \tcell{Appui aérien de l'Air Support Division.} \\
\altrow
\tkey{Bomb Squad} & \tcell{Unité de déminage. Déploiement sur menace explosive.} \\
\tkey{Richard} & \tcell{Unité de la Metropolitan Division. Présence sur les appels à risque.} \\
\altrow
\tkey{David} & \tcell{Unité d'intervention de la Metropolitan Division. Présence en code rouge et sur opération planifiée.} \\
\tkey{K9} & \tcell{Unité cynophile. Déploiement sur demande pour fouille de véhicule ou recherche de personne.} \\
\altrow
\tkey{William} & \tcell{Unité du Detective Bureau.} \\
\bottomrule
\end{tabular}

\cl[lincoln]{Les patrouilles en Lincoln sont interdites, sauf lorsque l'effectif disponible ne permet pas la constitution d'une Adam ou sur autorisation préalable du Command Staff ou de l'État-Major. Les Police Officer I probatoires ne peuvent patrouiller en Lincoln en aucune circonstance.}

\cl{Toute unité prend un indicatif à sa prise de service et le conserve pour l'ensemble de la garde. Le changement d'indicatif en cours de garde est annoncé à la radio et validé par le Watch Commander.}

% ---------------------------------------------------------------------
\article{Divisions et unités spécialisées}

\cl[divisions]{Le département comprend les divisions et unités suivantes.}

\vspace{4pt}
\noindent
\begin{tabular}{P{16mm} P{62mm} P{22mm} P{34mm}}
\headrow
\thead{Sigle} & \thead{Dénomination} & \thead{Type} & \thead{Grade minimal} \\
\tkey{D.B.} & \tcell{Detective Bureau} & \tcell{\textbf{Majeure}} & \tcell{Police Officer II} \\
\altrow
\tkey{Metro} & \tcell{Metropolitan Division} & \tcell{\textbf{Majeure}} & \tcell{Police Officer II} \\
\tkey{L.S.P.A.} & \tcell{Los Santos Police Academy} & \tcell{Mineure} & \tcell{Police Officer II} \\
\altrow
\tkey{M.R.D.} & \tcell{Management Relations Division} & \tcell{Mineure} & \tcell{Police Officer I} \\
\tkey{A.S.D.} & \tcell{Air Support Division} & \tcell{Mineure} & \tcell{Police Officer I} \\
\altrow
\tkey{K9} & \tcell{K9 Unit} & \tcell{Mineure} & \tcell{Police Officer I} \\
\tkey{T.D.} & \tcell{Traffic Division} & \tcell{Mineure} & \tcell{Police Officer I} \\
\bottomrule
\end{tabular}

\vspace{3pt}
{\lspdsans\fontsize{8}{10}\selectfont\color{lspdgrey}
Le SWAT ne constitue pas une division autonome : il s'agit d'une spécialité interne à la Metropolitan Division, accessible à ses seuls membres confirmés et habilités.}

\cl{Chaque division, unité et spécialité est soumise au présent règlement. Les protocoles internes de division le complètent et ne peuvent y déroger.}

\cl[refus-division]{Le Commanding Officer d'une division ne peut refuser une candidature sans motif. Sa décision est motivée par écrit. Le candidat éconduit peut former un recours devant le Command Staff. La nécessité de sélectionner les candidats au regard du nombre de candidatures constitue un motif de refus légitime.}

\cl{En cas de situation critique nécessitant le déploiement d'une unité spécialisée, le dispatcher peut l'autoriser sous l'accord du superviseur en charge.}

\sstitre{Organisation interne des divisions}

\begin{itemize}
  \item \stress{Commanding Officer} : Commanding Office, Director pour la Police Academy, Commander pour la K9 Unit.
  \item \stress{Supervisor} : encadrement intermédiaire de la division.
  \item \stress{Operator, Pilot ou Member} : membre confirmé, selon la division.
  \item \stress{Trainee} : membre en cours d'intégration, non autorisé à opérer seul aux couleurs de la division.
\end{itemize}

\cl[cumul]{Un agent peut appartenir à \textbf{une seule division majeure} et à \textbf{autant de divisions mineures qu'il le souhaite}, sous réserve d'assurer effectivement son service de patrouille.}

\cl{L'appartenance à une division n'emporte aucune autorité hiérarchique sur les agents extérieurs à celle-ci. Le grade de division ne prime jamais le grade de service.}
